% ============================================================================
%  Capítulo 8 — Conclusiones
% ============================================================================
\chapter{Conclusiones}
\label{cap:conclusiones}

\section{Conclusiones por objetivo}

\textbf{OE1 --- Caracterizar las perturbaciones.} Las perturbaciones del lazo
de oxígeno quedaron cuantificadas sobre 22 meses de operación
(Tabla~\ref{tab:ident_resumen}) y llevadas con \eqref{eq:pid_eps} a la unidad que decide el diseño,
\% de la señal de control. Con ellas se generó la perturbación sintética
(Figura~\ref{fig:ident_perturbacion}) con la que se validaron los controladores implementados, y en la
Figura~\ref{fig:pid_adrc_perturbacion} se ve que durante las 12~h el observador sigue la perturbación
aplicada. El $95~\%$ de su energía está por debajo de $12$~rad/h y el $99~\%$ por debajo de $56$~rad/h
(Tabla~\ref{tab:ident_espectro_f}), que es lo que orientó el ancho de banda del observador.

\textbf{OE2 --- Diseñar un GPI robusto sobre el error.} El controlador GPI de
orden $n=1$, $m=1$ quedó diseñado, discretizado y dotado de un esquema de
\emph{anti-windup} deducido de sus propias ganancias. El esquema elimina el sobrepaso por saturación:
en el arranque, sin él el oxígeno sube hasta $2{,}79$~mg/L, y con él llega a la consigna sin pasarse
(Figura~\ref{fig:sat_arranque}). Se verificó que su diseño en tiempo continuo coincide con su
realización en discreto (Figura~\ref{fig:gpi_discreto_vs_adrc}). Además, se implementó en el PLC para
los dos tipos de sintonía, por anchos de banda o por las ganancias $k_0$, $k_1$ y $k_2$,
correspondientes a las de un PI con filtro del error \eqref{eq:gpi_pi_filtro}.

\textbf{OE3 --- Diseñar un ADRC basado en ESO sobre el error.} El controlador
quedó diseñado con criterio físico para las dos perillas: $\wc = 4/\tau$, cuatro veces más rápido que
la planta en lazo abierto, y $\wo = 5\,\wc = 60$~rad/h, validado contra el contenido espectral de la
perturbación y con un rechazo coherente con las perturbaciones de orden $m=1$. Con esa sintonía el
margen de fase es de $65{,}8^{\circ}$ sin retardo y de $36{,}0^{\circ}$ con $\theta = 50$~s
(Figuras~\ref{fig:adrc_frecuencia} y \ref{fig:adrc_margenes_retardo}), y el escalón de $2{,}0$ a
$2{,}2$~mg/L se asienta en $15{,}7$~min. Además, se presenta una relación de equivalencia entre el GPI
y el ADRC basado en observador de estado extendido \eqref{eq:eso2gpi}.

\textbf{OE4 --- Metodología de implementación en PLC.} Se documentó una
metodología (Sección~\ref{sec:impl_resumen}), verificada sobre plataformas basadas en CODESYS, con
bloques de función en texto estructurado formulados sobre el error para que
sustituyan al PID existente sin modificar la arquitectura del lazo
(Tablas~\ref{tab:plc_adrc_eso}, \ref{tab:plc_gpi} y \ref{tab:plc_adrc_gpi}, Figura~\ref{fig:plc_codesys}, Anexo~\ref{anx:st}). Los dos
controladores se ejecutaron en las plataformas, lo que valida su funcionamiento junto con los
productos digitales diseñados bajo los aspectos prácticos de implementación, como la curva de Bézier
para la referencia y el prefiltro (Tablas~\ref{tab:plc_prefilter_gpi}, \ref{tab:plc_prefilter_adrc} y \ref{tab:plc_bezier}).

\textbf{OE5 --- Validación frente a un PID convencional.} La validación se hizo en simulación, con la
planta identificada y la perturbación sintetizada de los datos reales, ya que el controlador no se ha
puesto en marcha en campo (Sección~\ref{sec:limitaciones}). La comparación se hizo contra dos
referencias: los arranques medidos del PID real, que unas veces llegan limpios y otras oscilan con
sobreimpulsos de $44$ y $91~\%$ (Figura~\ref{fig:ident_pid}), y un PID ISA sintonizado para que su
señal de control no salte más que la del ADRC en el cambio de consigna y tenga la misma sensibilidad
máxima sin retardo. Con esas condiciones, el ADRC es mejor en el escalón de consigna, con un IAE de
$42{,}7$ frente a $88{,}4$~mg$\cdot$s/L con el retardo de diseño (Tabla~\ref{tab:pid_escalon},
Figura~\ref{fig:pid_adrc_escalon}), y en el rechazo de una perturbación de entrada, que corrige sin
pasarse de la consigna (Tabla~\ref{tab:pid_pert_escalon}). En las 12~h con perturbación sintética el
IAE del ADRC es algo menor, $1{,}86$ frente a $2{,}23$~mg$\cdot$h/L (Sección~\ref{sec:pid_perturbacion}), y
el PID se pasa unas tres veces más en los cambios de consigna y maneja peor el ruido, con una señal de
control más ruidosa ($\sigma_u$ de $0{,}22$ frente a $0{,}009~\%$,
Figura~\ref{fig:pid_adrc_escalon}). El PID solo queda mejor en robustez con retardo, con
$M_s$ de $1{,}59$ frente a $2{,}32$ (Tabla~\ref{tab:pid_robustez}), a cambio de poco más de un tercio
del ancho de banda. Para esta planta de primer orden, un PI con el mismo margen de fase del ADRC puede
superarlo en rechazo (Figura~\ref{fig:pid_mejores}), pero deja pasar mucho más ruido al soplador.

\section{Conclusión general}

Las dos arquitecturas de rechazo activo evaluadas --- GPI y ADRC con ESO ---
son, para este proceso y con las ganancias correspondientes, el mismo
controlador escrito de dos maneras, es decir, son equivalentes. Su ventaja frente al PID no está
estrictamente en el IAE, dado que con diferentes sintonizaciones puedo tener mejores IAE. Bajo la misma
sensibilidad máxima y el mismo esfuerzo inicial, está en el precio que cada uno paga por ese IAE: en
sobrepico, en el rechazo de perturbaciones y en el ruido que llega al soplador. En general, se pudieron validar los controladores y sus respectivas
equivalencias en una problemática industrial donde el PID es la forma habitual de control, y se llegó
a una metodología y a unos aspectos prácticos de implementación para realizar una puesta en marcha
adecuada a cada problemática que pueda llegar a surgir, usando plataformas basadas en CODESYS. Lo que
queda pendiente es confirmarlo en campo.

\section{Aportes principales}

\begin{enumerate}

  \item \textbf{Validación de controladores más que todo teóricos, pero con enfoque industrial.}
  Si bien ya se ha demostrado que estos tipos de controladores avanzados, como el ADRC y el GPI, tienen
  algunas ventajas frente a controladores clásicos como el PID, a nivel de manejo de ruido y de
  facilidad de sintonización relacionada con fenómenos físicos como los anchos de banda de la planta y
  de la perturbación, o incluso por \emph{pole placement}, no son tan vistos en proyectos o procesos
  industriales, donde el PID sigue ganando terreno. Aquí se demuestra una forma más práctica de
  implementarlos, diseñados sobre los datos reales de una planta en operación.

  \item \textbf{La reconstrucción explícita de la equivalencia ESO--GPI.} Se buscaba una forma de
  obtener el mismo resultado bajo dos enfoques diferentes de construcción de controladores, y la
  equivalencia \eqref{eq:eso2gpi} se mantiene en tiempo continuo, en discreto y también con saturación
  (Figuras~\ref{fig:gpi_vs_adrc}, \ref{fig:sat_arranque} y \ref{fig:gpi_discreto_vs_adrc}).

  \item \textbf{Una metodología de implementación en PLC verificada en CODESYS.} Se entregan códigos
  escritos en texto estructurado, correspondiente a la norma IEC~61131-3 (Anexo~\ref{anx:st}), y su
  metodología de sintonización y uso (Sección~\ref{sec:impl_resumen}), por lo que queda listo para que
  cualquier usuario lo pruebe y tenga forma de sintonizarlo fácilmente.

  \item \textbf{Aspectos prácticos de implementación.} Se realizaron diferentes experimentos para
  demarcar qué se debe considerar a la hora de integrar estos tipos de controladores, o en general
  controladores, en entornos industriales, y cómo se puede hacer un mejor uso de ellos con
  acondicionamientos específicos como el manejo de la saturación, una referencia acotada e incluso un
  manejo del ruido, que a nivel industrial puede estar presente en todos los procesos.

\end{enumerate}

\section{Trabajo futuro}

\begin{enumerate}
  \item \textbf{Terminar la puesta en marcha en campo.} Teniendo ya todos los esquemas listos para
  probar, simulados y con sus códigos establecidos, queda su puesta en marcha como pendiente.

  \item \textbf{Evaluar bajo influentes de lluvia, cambios ambientales y zonas muertas de los motores.}
  Si bien ya está validado bajo perturbaciones y sintonización relacionadas con la planta, falta
  validarlo bajo históricos de un tiempo prolongado, de meses, para ver su comportamiento bajo
  condiciones ambientales cambiantes, como tormentas o lluvias, donde se modifica el caudal y cambia la
  biología de la planta. Si bien el controlador responde adecuadamente a perturbaciones supuestas e
  interpretadas de los datos, es necesario realizar pruebas prácticas en las que se involucren mayores
  perturbaciones, y más que todo dirigidas a los actuadores, como cambios en los difusores o en la
  potencia de los mismos motores al cambiar equipos o introducir nuevos, siendo esto evidenciado en un
  cambio de las ganancias, además de la presencia de posibles zonas muertas en los motores.

  \item \textbf{Reducción del esfuerzo de control.} Una de las diferencias entre el esquema basado en el
  error y el esquema de estimación de la señal de salida es que, al introducir la referencia por el
  observador en el error de seguimiento y meter un paso, la acción de control va a ser muy fuerte,
  debido a que el controlador debe compensar la derivada de un paso, que es un impulso,
  $u^* = [y^*]^{(n)}$, y eso genera un enorme esfuerzo de la señal de control. Aquí se redujo con Bézier y
  con el prefiltro; otra forma de reducirlo es mediante un \emph{clutch} que no inyecte directamente todo
  el esfuerzo de la señal de control a la planta, sobre lo que ya existen documentos como
  \cite{quecan2026clutch}.

  \item \textbf{Implementación de controladores ADRC con otro tipo de perturbación.} Poder llevar a
  estructuras implementables en equipos digitales grados mayores de perturbación, para tener mejor
  seguimiento ante otros tipos de señales, como rampas o incluso señales oscilatorias, de las que ya
  existe una teoría grande, y el único esfuerzo es ver la forma de facilitar su implementación.

  \item \textbf{Implementación analógica de controladores en circuitos dedicados.} Si bien la
  discretización ha sido una forma muy buena de implementar diferentes controladores teóricos continuos,
  sería interesante validar estos tipos de controladores en SoC con periféricos totalmente analógicos,
  cuya única función sea realizar el control. Por lo que, bajo un enfoque de diseño ASIC
  (\emph{Application-Specific Integrated Circuit}) mixto, se podría llegar a crear un producto con una
  parte digital que permita la interacción con el usuario y modelar cada una de las ganancias del
  controlador por software, pero con una realización completamente analógica en un chip dedicado.

\end{enumerate}
